\documentclass[11pt]{article}
\usepackage[margin=1in]{geometry}
\usepackage[T1]{fontenc}
\usepackage{lmodern,amsmath,amssymb,amsthm,booktabs,microtype}
\usepackage[hidelinks]{hyperref}
\newtheorem{theorem}{Theorem}
\newtheorem{lemma}[theorem]{Lemma}
\newtheorem{proposition}[theorem]{Proposition}
\newcommand{\Fp}{\mathbb F_p}
\newcommand{\Zp}{\mathbb Z_{(p)}}
\newcommand{\fl}[1]{\left\lfloor #1\right\rfloor}
\newcommand{\ff}[2]{#1^{\underline{#2}}}
\newcommand{\val}{v_p}
\newcommand{\res}{\operatorname*{Res}}
\hypersetup{pdftitle={The irrationality exponent of π² is less than 4.589897},
            pdfauthor={Ryan Matthew Casper}}
\title{The irrationality exponent of \texorpdfstring{$\pi^2$}{pi squared}\\is less than 4.589897}
\author{Ryan Matthew Casper}
\date{November 25, 2026}
\begin{document}
\maketitle
\begin{abstract}
We prove that the irrationality exponent of $\pi^2$ is less than
$4.589897$. Leibtag, Barkan, Kalisch, Weinbaum, Shalyt, and Kaminer
obtained the bound $4.62675$ using rational approximations to
$\zeta(2)$. They removed a common prime factor $p$ from the numerator
and denominator whenever $p^2>50n$ and $n/p$ lay in one of six
intervals. Under the same cutoff, we show that every positive integer
translate gives primes that divide both coefficients even after this
normalization. The product of these additional
primes has logarithm $(0.1086823692\ldots+o(1))n$. Dividing by this
product reduces both the denominator and the error of the associated
integer linear form, giving the stated bound.
\end{abstract}

\section{The bound and the primes}
Since $\pi$ is transcendental, so is $\pi^2$. We study how closely
rational numbers can approximate it. For an irrational real number
$\xi$, let $\mu(\xi)$ denote the supremum
of the real $\nu$ for which $|\xi-a/b|<b^{-\nu}$ has infinitely many
solutions in reduced fractions $a/b$, $b>0$. Multiplication by a nonzero
rational preserves this exponent, so Euler's identity
$\pi^2=6\zeta(2)$ lets us work with $\xi=\zeta(2)$.

\begin{theorem}\label{th:main}
The irrationality exponent of $\pi^2$ satisfies $\mu(\pi^2)<4.589897$.
\end{theorem}

In particular, the inequality $|\pi^2-a/b|<b^{-4.589897}$ holds
for only finitely many reduced fractions $a/b$ with $b>0$.

Our earlier theorem on quadratic periods gives $\mu(\pi/\sqrt d)=2$
for each positive squarefree integer $d$ \cite[Theorem~1(1)]{Casper}.
It gives a stronger estimate when the approximants lie in a fixed
squareclass. Indeed, write $a/b=d(u/v)^2$ with both fractions positive
and reduced. Then $b=v^2/\gcd(d,v^2)\ge v^2/d$.
If $|\pi^2-a/b|<b^{-\nu}$ for a fixed $\nu>1$, taking square roots gives
$|\pi/\sqrt d-u/v|\ll_{d,\nu}v^{-2\nu}$. Since $2\nu>2$, only finitely
many such fractions exist for each fixed $d$. Theorem~\ref{th:main}
gives one bound for rational approximants in all squareclasses.

Zudilin obtained the bound $5.095412$ by comparing two hypergeometric
representations of the same linear forms \cite{Z}. Niedbala Giraudin
obtained $5.0193784$ by changing their parameters \cite{N}.
Leibtag et al.\ constructed a trajectory in a conservative matrix field
and proved $\mu(\zeta(2))<4.62675$ \cite{L}. We use their first-column
approximants and their estimates for growth, decay, and independence.
The improvement comes from dividing both integer coefficients by further
primes.\footnote{Carella proposed $\mu(\pi^2)=2$ in
\cite[Theorem 1.1]{C}. For the convergents $a_j/b_j$, put
$z_j=|\pi^2 b_j-a_j|$. Equation (6) gives $|\sin z_j|\asymp z_j$,
while (8) also needs $b_j^{-1}\ll |\sin z_j|$. If $\mu(\pi^2)>2$,
Legendre's criterion gives a subsequence with $z_j<b_j^{-1-\delta}$
for some $\delta>0$. Then $b_j|\sin z_j|\le b_jz_j\to0$ along this
subsequence, so (6) alone does not give the lower bound in (8).}

Write $\mathcal I$ for the collection of intervals
\begin{equation}\label{eq:intervals}
 \left(\tfrac18,\tfrac2{13}\right),\
 \left(\tfrac29,\tfrac3{13}\right),\
 \left(\tfrac3{10},\tfrac4{13}\right),\
 \left(\tfrac38,\tfrac5{13}\right),\
 \left(\tfrac49,\tfrac6{13}\right),\
 \left(\tfrac89,\tfrac{12}{13}\right),
\end{equation}
and let $I$ be their union. The first-column normalization in \cite{L}
already removes one factor $p$ when $n/p\in I$ and $p^2>50n$.
Let $(P_n,Q_n)$ be the resulting integral pair, given explicitly in
\eqref{eq:pq}. The following theorem gives a further common divisor.

\begin{theorem}\label{th:div}
If $n$ is a positive integer, $p$ is prime, $50n<p^2$, and
$n/p\in t+I$ for an integer $t\ge1$,
then $p$ divides both $P_n$ and $Q_n$.
\end{theorem}

The original six intervals use primes larger than $n$. Every translated
interval uses primes smaller than $n$. To see why the same intervals
occur again, write $n=tp+r$ and group the coefficient indices as
$j=k+p\ell$, with $0\le k<p$. Divide each coefficient by the common
power of $p$ coming from its carries and reduce modulo $p$. It then
splits into two factors. The low factor depends on $k$ and the
remainder $r$. The high factor can change with $t$ and $\ell$, but
is constant as $k$ varies over the block. The low factors sum to zero in
$\Fp$. For the numerator we will need this cancellation twice, once
with harmonic weights and once in a finite difference.

\section{The first-column forms}
We recall the finite first-column formula in \cite[\S6 and Appendix B]{L}
and its detailed derivation in \cite{LC}. Write $m=11n$ and put
\[
 F(s)=\prod_{u=9}^{11}\binom{s+un}{8n},\qquad
 c_j=(-1)^j\binom mj F(j),\qquad
 R(s)=\frac{(-1)^m m!\,F(s)}{\prod_{j=0}^{m}(s-j)} .
\]
Unless otherwise specified, products or sums indexed by $u$ range
over $9,10,11$. The polynomial part $S$ of $R$ satisfies
\begin{equation}\label{eq:partial}
 S(s)=R(s)-\sum_{j=0}^m\frac{c_j}{s-j},\qquad \deg S\le13n-1 .
\end{equation}
With $\Delta f(s)=f(s+1)-f(s)$ and $H_b^{(2)}=\sum_{i=1}^b i^{-2}$,
define
\begin{align}
 B_n&=\sum_{j=0}^m c_j,\qquad
 J_\eta(S)=\sum_{i=0}^{13n-1}\frac{(-1)^i}{i+1}\Delta^i S(-\eta),
 \nonumber\\
 A_n&=\sum_{j=0}^m c_j H_{j+\eta}^{(2)}+J_\eta(S).
 \label{eq:finite}
\end{align}
The finite first-column identity of \cite{LC} allows any integer
$2n\le\eta\le9n$. We will choose a multiple of $p$ so that, after
scaling by $p^2$, the harmonic weight is constant modulo $p$ on each block. To see the invariance directly, note that
$J_{\eta+1}(S)-J_\eta(S)=-S'(-\eta-1)$, as follows by testing
$S(s)=e^{zs}$ coefficient by coefficient. The sum defining $J_\eta$
on this series agrees with $e^{-\eta z}z/(e^z-1)$ through degree
$13n-1$ in $z$, enough for every polynomial under consideration.
Thus the change in $A_n$
is $-R'(-\eta-1)$ by \eqref{eq:partial}. At
$-9n\le s\le-2n-1$, the binomials with $u=9,10$ both vanish, so
$R'(s)=0$.

Write $L_j=\operatorname{lcm}(1,\ldots,j)$. For $0\le y<1$ set
\[
 \rho(x,y)=[y>\{11x\}]+\sum_u[\{ux+y\}<\{8x\}],\qquad
 \phi(x)=\min_{0\le y<1}\rho(x,y).
\]
Here brackets denote indicators and braces denote fractional parts.
This is the carry function of \cite{L}, written with fractional
parts. The normalization uses two divisors,
\[
 G_n=\prod_{\sqrt{50n}<p\le13n}p^{\phi(n/p)},\qquad
 E_n=\prod_{\substack{p^2>50n\\ n/p\in I}}p ,
\]
where both products run over primes. With $C_n=F(0)$, the
first-column entries $\mathsf P_{n,0},\mathsf Q_{n,0}$ in the finite
identity of \cite{LC} satisfy
$C_n(\mathsf P_{n,0},\mathsf Q_{n,0})=2(A_n,B_n)$.
The integral normalization in \cite[Lemmas 6.20--6.21]{L} is thus
\begin{equation}\label{eq:pq}
 (P_n,Q_n)=2\,\frac{L_{12n}L_{13n}}{G_nE_n}\,(A_n,B_n)\ \in\mathbb Z^2 .
\end{equation}

\section{The low digit}
The companion proof of Leibtag et al.\ \cite{LC} constructs the
low polynomial (Q) and the kernel congruences (R), (C$'$), and
(S$'$) for the six original intervals. We show how their table
(T) changes when $n$ increases by a multiple of $p$, then derive the residue and interpolation identities for $p<n$.
Fix $n,p$ as in Theorem~\ref{th:div}. Thus $p<n$ and $p>50$.
Write $8n=ap+b$, $11n=Mp+d$, with $0\le b,d<p$, and let
$H=\phi(n/p)$. For each $0\le k<p$ put
\[
 k+un=p\lambda_u+r_u,\quad 0\le r_u<p,\qquad
 \delta_u=[r_u<b],\quad e=[k>d],\quad \chi(k)=e+\sum_u\delta_u.
\]
Then $H=\min_k\chi(k)$. In fact, between $k/p$ and $(k+1)/p$
all the translated fractional-part tests in $\rho(n/p,y)$ equal
their values at $k/p$, apart from the first strict test, which can
only increase. In particular the discrete minimum is attained.

\begin{lemma}\label{lem:table}
On the set $\chi(k)=H$ the tuple $(e,\lambda_9,\lambda_{10},
\lambda_{11},\delta_9,\delta_{10},\delta_{11})$ is independent of $k$.
Also $13n<(3a-M+H)p$.
\end{lemma}
\begin{proof}
Take $x=r/p$, $r=n-tp$. The table below is the first-column table of \cite{LC} evaluated
at $r$. In this table only, $a=\lfloor8r/p\rfloor$,
$M=\lfloor11r/p\rfloor$, $\lambda_u=\lfloor(k+ur)/p\rfloor$,
and $K=3a-M+H$.
\[
\begin{array}{c|rr|rr|c|c|r}
 x&a&M&H&e&(\lambda_9,\lambda_{10},\lambda_{11})&
      (\delta_9,\delta_{10},\delta_{11})&K\\ \midrule
 (1/8,2/13)&1&1&0&0&(1,1,1)&(0,0,0)&2\\
 (2/9,3/13)&1&2&2&0&(2,2,2)&(1,1,0)&3\\
 (3/10,4/13)&2&3&1&0&(2,3,3)&(0,1,0)&4\\
 (3/8,5/13)&3&4&0&0&(3,3,4)&(0,0,0)&5\\
 (4/9,5/11)&3&4&1&0&(4,5,5)&(0,1,0)&6\\
 (5/11,6/13)&3&5&2&1&(4,5,5)&(0,1,0)&6\\
 (8/9,10/11)&7&9&0&0&(8,9,10)&(0,0,0)&12\\
 (10/11,12/13)&7&10&1&1&(8,9,10)&(0,0,0)&12
\end{array}
\]
To recover the table, split $0\le y<1$ at the seven fractional parts
\[
 \{11x\},\quad \{-ux\},\quad \{(8-u)x\}\qquad (u=9,10,11).
\]
The first cut tests $e$, the next three the jumps of $\lambda_u$,
and the last three mark
where $\{ux+y\}$ reaches $\{8x\}$. Test the cuts themselves as well
as the open pieces, using the strict comparisons in $\rho$.
Writing the cuts as $\{cx\}$, their order changes only where $cx$
or $(c-c')x$ is integral.
These are the base-period cells in \cite{LC}. The two internal
endpoints with denominator $11$ would force $p=11$.
In each row $13x<K$. Replacing $r$ by $n=r+tp$ keeps
$b,d,r_u,e,\delta_u,\chi(k)$ fixed for every $k$. It increases
$\lambda_u$ by $ut$, $a$ by $8t$, $M$ by $11t$, and $K$ by $13t$.
For example, the first row gives $13r<2p$. After $t$
translations this becomes $13(r+tp)<(2+13t)p$. The same increase
on both sides preserves the strict inequality in every row.
\end{proof}

Over $\Fp$ consider the rational expression
\begin{equation}\label{eq:low}
 U(X)=\frac{\prod_u\binom{X+un}{b}}
 {(X^p-X)^{H-1}\prod_{j=0}^d(X-j)} .
\end{equation}
At $k$, a binomial in the numerator has a simple zero exactly when
$r_u<b$. Hence $U$ is a polynomial with multiplicity $\chi(k)-H$
at $k$. This also covers $H=0$, where $X^p-X$ occurs upstairs.
Its degree is $3b+p-1-d-pH$. Since $3b-d=13n-(3a-M)p$,
Lemma~\ref{lem:table} gives $3b<d+pH$, so the degree is below $p-1$. The sum of every power $X^i$, $0\le i<p-1$,
over $\Fp$ is zero in $\Fp$. We therefore have
\begin{equation}\label{eq:zerosum}
 U(k)\ne0\ \Longleftrightarrow\ \chi(k)=H,\qquad
 \sum_{k=0}^{p-1}U(k)=0 .
\end{equation}

\section{Residues and interpolation}
Use the tuple from Lemma~\ref{lem:table} to form a rational function
over $\Fp$ in the high digit $V$. With
$\ff{x}{h}=x(x-1)\cdots(x-h+1)$, put
\[
 G(V)=\frac{\displaystyle\prod_u
          \frac{\ff{(V+\lambda_u)}{a+\delta_u}}{a!}}
                {\displaystyle\prod_{\ell=0}^{M-e}(V-\ell)},\qquad
 C=(-1)^{d-M+H-1}M!\,d! .
\]
Write $r_\ell=\res_{V=\ell}G(V)$ for $0\le\ell\le M$, taking
residue zero at a regular point. The poles in this formula are
distinct since $M<p$. We take congruences in $\Zp$, the rationals whose reduced
denominator is prime to $p$.

\begin{lemma}\label{lem:factor}
Extend $c_j$ by zero beyond $m$. For $0\le k<p$ and $0\le\ell\le M$,
\begin{equation}\label{eq:coeff}
 p^{-H}c_{k+p\ell}\equiv C U(k)r_\ell \pmod p .
\end{equation}
For integers $M+1\le V\le p-M-2$,
\begin{equation}\label{eq:away}
 p^{1-H}R(k+pV)\equiv C U(k)G(V) \pmod p .
\end{equation}
\end{lemma}
\begin{proof}
For $0\le vp+w<p^2$, $0\le w<p$, Wilson's theorem in each complete
block gives $p^{-v}(vp+w)!\equiv(-1)^v v!\,w!$. We apply this formula
first to $c_j$ at $j=k+p\ell\le m$. Its valuation is
$e+\sum_u\delta_u=\chi(k)$: these are the borrows in $m-j$ and in the
three subtractions $j+un-8n$. The largest factorial argument is at
most $22n<p^2$. Suppose first that $\chi(k)=H$. The low factorials
give the unit
\[
 L(k)=\frac{1}{k!(d-k+ep)!}
       \prod_u\frac{r_u!}{b!(r_u-b+\delta_up)!}.
\]
At $k$ remove all factors $X-k$ before evaluating \eqref{eq:low}.
The corresponding unit of $X^p-X$ is $-1$. Each shifted falling
factorial of length $b$ contributes
$(-1)^{\delta_u}r_u!/(r_u-b+\delta_up)!$.
The inverse unit of $\prod_{j=0}^d(X-j)$ is
$(-1)^{d-k}/[k!(d-k+ep)!]$. These evaluations follow from
$z!(p-1-z)!\equiv(-1)^{z+1}$. Combining the signs gives
\[
 U(k)=(-1)^{1-e+d-k}L(k).
\]
Restoring the high factorials in $c_j$ yields
\[
 p^{-H}c_j\equiv
 (-1)^{k+\ell+H}d!L(k)\frac{M!}{\ell!(M-\ell-e)!}
      \prod_u\frac{\ff{(\ell+\lambda_u)}{a+\delta_u}}{a!}.
\]
The derivative of $\prod_{v=0}^{M-e}(V-v)$ at $\ell$ is
$(-1)^{M-e-\ell}\ell!(M-e-\ell)!$. This proves \eqref{eq:coeff}
at each included minimum. For larger carry, both sides vanish.

In the final block $\ell=M$ we omit $k>d$. If the minimum tuple
has $e=0$, those omitted residues have $U(k)=0$. If it has $e=1$,
$G$ is regular at $M$, so $r_M=0$. Thus \eqref{eq:coeff} holds on
the completed final block as well.

For \eqref{eq:away}, write $s=k+pV$ and use
$R(s)=(-1)^m m!F(s)(s-m-1)!/s!$. Now $s>m$ and $s+m<p^2$.
The high digit of $s-m-1$ is $V-M-1+e$, and its low digit is
$p-1-(d-k+ep)$. Counting the high digits gives
$\val(R(s))=\chi(k)-1$. At a minimum the high factorial ratio
$(V-M-1+e)!/V!$ supplies the denominator of $G(V)$, and the three
binomials supply its numerator. For the low factorial left by
$(s-m-1)!$, use the complement formula above. The signs from
$(-1)^m$, the factorial units, and the complement are respectively
$(-1)^{M+d}$, $(-1)^{H-1}$, and $(-1)^{d-k+e+1}$.
Together with the formula for $U(k)$, this gives $C U(k)G(V)$.
Higher carry again gives zero.
\end{proof}

Congruence \eqref{eq:away} gives $R$ away from the poles of $G$.
The numerator in \eqref{eq:finite} needs the polynomial part $S$ at
negative integers, including classes where $G$ has a pole. At nodes
away from the poles, both $R$ and its pole terms have the factor
$U(k)$ after scaling and reduction. We subtract those terms there,
then interpolate the remaining high-digit factor. The same polynomial
$J$ will serve for every $k$.

\begin{lemma}\label{lem:interpolate}
There is a polynomial $J\in\Fp[V]$ such that, for all integers $v$
and all $0\le k<p$,
\[
 T(k+pv)\in\Zp,\qquad T(k+pv)\equiv U(k)J(v)\pmod p,\qquad
 T=p^{1-H}S .
\]
\end{lemma}
\begin{proof}
First $T$ takes values in $\Zp$ at every integer.
To see this, take the $13n$ consecutive nodes $s=-13n,\ldots,-1$.
Every nonzero factor of $F(s)$ and of the pole denominator in $R(s)$
has absolute value at most $24n<p^2$. Write $s=k+pv$ with $0\le k<p$.
In the $u$th falling factorial of $F(s)$ there are $a+\delta_u$
multiples of $p$, compared to $a$ in $(8n)!$.
The pole denominator has $M+1-e$ such factors, compared to $M$ in $m!$.
Hence $R(s)=0$ or $\val(R(s))=\chi(k)-1$.
Every $c_j/(s-j)$ has valuation at least $H-1$ by the same bound
on $|s-j|$ and the coefficient valuation in Lemma~\ref{lem:factor}.
Equation~\eqref{eq:partial} gives $T(s)\in\Zp$ at these nodes.
Newton interpolation in the basis $\binom{s+13n}{i}$, using
$\deg T\le13n-1$, gives this integrality
at every integer and hence $\Delta^iT(0)\in\Zp$ for every $i$.

Set $D=\lfloor(13n-1)/p\rfloor$. For each fixed $k$ the reductions
of $T(k+pv)$ are the values of a polynomial $W_k(v)$ over $\Fp$
of degree at most $D$. Indeed, expand $T$ in
$\binom{s}{i}$ with coefficients $\Delta^iT(0)$. If $i=hp+z$,
$0\le z<p$, then
\[
 \binom{k+pv}{hp+z}\equiv\binom kz\binom vh\pmod p .
\]
This identity holds even for negative integers $v$: compare
coefficients in $(1+X)^{k+pv}\equiv(1+X)^k(1+X^p)^v$ as formal
power series. Here $h\le D<p$, so $\binom vh$ has degree $h$ over
$\Fp$.

Take the $D+1$ nodes $V=M+1,\ldots,M+D+1$. They lie in the range
of \eqref{eq:away}, since $2M+D+3\le35n/p+3<p$.
In \eqref{eq:partial} the terms with $j\not\equiv k\bmod p$ vanish
after multiplying by $p^{1-H}$ and reducing. For $j=k+p\ell$ the
reduced term is $C U(k)r_\ell/(V-\ell)$, including the zero extension
of the final block. It follows that
\[
 W_k(V)=U(k)\, C\left(G(V)-\sum_{\ell=0}^M\frac{r_\ell}{V-\ell}\right)
\]
at each of these nodes. Interpolate the $D+1$ bracketed values
multiplied by $C$ to obtain a single polynomial $J$ of degree at most
$D$. Then $W_k=U(k)J$ for every $k$, since both sides have degree
at most $D$ and agree at $D+1$ distinct points.
\end{proof}

\section{A common divisor of the two coefficients}
\begin{proof}[Proof of Theorem~\ref{th:div}]
Write $H=\phi(n/p)$. We first show that $v_p(B_n)\ge H+1$
and $v_p(A_n)\ge H-1$. We then combine these bounds with the valuation of the
multiplier in \eqref{eq:pq}.

The denominator estimate follows from \eqref{eq:coeff} and
\eqref{eq:zerosum}: in each completed block $\sum_k c_{k+p\ell}$
is divisible by $p^{H+1}$. To make the harmonic weight constant on
each block, choose
$\eta=p\lceil2n/p\rceil=pq$. Then $2n\le\eta<3n$, so \eqref{eq:finite}
applies. For $j=k+p\ell\le m$ we have $j+\eta<14n<p^2$, whence
\[
 p^2 H_{j+\eta}^{(2)}\equiv\sum_{w=1}^{\ell+q}w^{-2}\pmod p .
\]
This weight is constant on the block. Thus \eqref{eq:coeff} and
\eqref{eq:zerosum} show that $p^{2-H}$ times the harmonic sum vanishes
modulo $p$, giving valuation at least $H-1$.

It remains to show $J_\eta(T)\in\Zp$, with $T$ as in
Lemma~\ref{lem:interpolate}. All its finite differences are integral
at $-\eta$. Since $i+1\le13n<p^2$, the only weights $1/(i+1)$
that can cost a factor of $p$ occur at $i=vp-1$. Write the index in
the finite difference as $up+k$, with $0\le u<v$ and $0\le k<p$.
Lucas's congruence and $\binom{p-1}{k}\equiv(-1)^k$ give
\[
 (-1)^{vp-1-up-k}\binom{vp-1}{up+k}
       \equiv(-1)^{v-1-u}\binom{v-1}{u}\pmod p .
\]
By Lemma~\ref{lem:interpolate}, the corresponding polynomial value
is $T(-\eta+up+k)\equiv U(k)J(u-q)$. Summing over $k$ gives zero
for each $u$, by \eqref{eq:zerosum}. Therefore
$p\mid\Delta^{vp-1}T(-\eta)$ in $\Zp$, as required. This proves
$\val(A_n)\ge H-1$.

Both least common multiples in \eqref{eq:pq} contain $p$ exactly once since
$p<n<13n<p^2$. The carry divisor contains $p^H$, and $E_n$ contains
only primes larger than $n$. The multiplier in \eqref{eq:pq} thus
has valuation $2-H$. Applying it proves the theorem.
\end{proof}

\section{The logarithmic saving}
Let $T_n$ be the product of the distinct primes in Theorem~\ref{th:div}.
For fixed $t\ge1$ and $(\alpha,\beta)\in\mathcal I$, the primes in
question lie in
\[
  \frac{n}{t+\beta}<p<\frac{n}{t+\alpha}.
\]
For fixed $t$ these inequalities imply $p^2>50n$ for all large $n$.
The intervals for different $t$ and different components of $I$ are
disjoint. Hence the prime number theorem, in the form
$\vartheta(x)=\sum_{p\le x}\log p\sim x$, gives rate
\[
 s_K=\sum_{(\alpha,\beta)\in\mathcal I}\sum_{t=1}^{K}
       \left(\frac{1}{t+\alpha}-\frac{1}{t+\beta}\right)
\]
for the product from $1\le t\le K$ (rate means logarithm divided by
$n$ in the limit). Endpoint contributions are $O_K(\log n)$.
All primes from $t>K$ are at most $n/K$, so their total logarithm is
at most $\vartheta(n/K)$. First letting $n$ tend to infinity, then
$K$, gives the exact rate
\begin{equation}\label{eq:rate}
 \lim_{n\to\infty}\frac{\log T_n}{n}
 =s=\sum_{(\alpha,\beta)\in\mathcal I}\sum_{t\ge1}
       \frac{\beta-\alpha}{(t+\alpha)(t+\beta)} .
\end{equation}
The infinite sum gives the rate quoted in the abstract. For the
irrationality exponent it suffices to keep the first million translates.

\begin{lemma}\label{lem:num}
For $K=10^6$ one has $s_K>s_0^*$, where $s_0^*=10868226/10^8$.
\end{lemma}
\begin{proof}
Write $(x)_j=x(x+1)\cdots(x+j-1)$, with $(x)_0=1$. Subtracting
successive remainders gives
\[
 \frac1{t+a}=\sum_{j=0}^J\frac{(1-a)_j}{(t+1)_{j+1}}
       +\frac{(1-a)_{J+1}}{(t+a)(t+1)_{J+1}} .
\]
If $0<a<b<1$, the last term is strictly larger at $a$ than at $b$.
Keep the first twelve periods exactly and drop the difference of
remainders thereafter, with $J=12$. Telescoping in $t$ gives the
following lower bound for the contribution of $(a,b)$:
\begin{align*}
 D(a,b)={}&\sum_{t=1}^{12}\frac{b-a}{(t+a)(t+b)}\\
   &+\sum_{j=1}^{12}\frac{(1-a)_j-(1-b)_j}{j}
       \left(\frac1{(14)_j}-\frac1{(K+2)_j}\right).
\end{align*}
The integer parts of $10^{12}D(a,b)$ in the
order of \eqref{eq:intervals} are given below. They can be obtained
by starting all rising products at $1$ and updating
$(z)_{j+1}=(z+j)(z)_j$.
\[
 \begin{array}{rrr}
 39328986343,&10507109706,&8689742248,\\
 10054938596,&16673511522,&23427973106 .
 \end{array}
\]
Their sum is $108682261521>10^{12}s_0^*$.
\end{proof}

The same expansion gives the decimal value in the abstract. Keep
the first $32$ periods, set $J=32$, and sum the rest to infinity.
For each interval $(a,b)$ the error in dropping the difference of
remainders is positive and at most
\[
 \sum_{t=33}^\infty\frac{(1-a)_{33}}{(t)_{34}}
 =\frac{(1-a)_{33}}{33(33)_{33}} .
\]
With the lower approximation obtained by dropping the remainders and
the upper one obtained by adding these six errors, the integer parts
after multiplication by $10^{14}$ are both $10868236921910$.

\section{Height and irrationality exponent}
We use the estimates (A7)--(A10) and the final arithmetic enclosure
in \cite{LC}. Before normalization, (A7)--(A8) give upper bounds of
$38.168733$ and $-24.467842$ for the logarithmic growth rates of
$|2B_n|$ and $|\xi(2B_n)-2A_n|$, respectively.
The multiplier $L_{12n}L_{13n}/(G_nE_n)$ adds its rate to both.
To bound its contribution, write
\[
 d_*=\frac{143297092558247390159}{13110552000000000000}.
\]
The section on the exponential cost of normalization in \cite{LC}
gives the rate $\gamma=\int_{1/13}^{\infty}\phi(x)x^{-2}\,dx$ for
$G_n$. Summing the six original prime intervals gives $233/120$ for
the rate of $E_n$, and the two least common multiples have combined
rate $25$. The rational enclosure for $\gamma$ in \cite{LC} therefore
bounds the rate of the multiplier by $25-\gamma-233/120\le d_*$.

Finally, dividing both coefficients by $T_n$ subtracts $s$ from
both logarithmic rates. Put $h=38.168733+d_*$ and
$g=38.168733+24.467842$. The resulting integer pairs
$(P'_n,Q'_n)$ satisfy
\[
 |Q'_n|\le e^{(h-s+o(1))n},\qquad
 |\xi Q'_n-P'_n|\le e^{-(g-h+s-o(1))n}.
\]
For all large $n$, the three successive first-column directions span
the three-dimensional state space. Indeed, if $A_n^{\rm tr}$ denotes
the normalized transition matrix of \cite{LC}, their matrix in the
level-$n$ basis is
$(e_0,A_n^{\rm tr}e_0,A_n^{\rm tr}A_{n+1}^{\rm tr}e_0)$, whose
limiting determinant is nonzero by (A9). The two readout rows are
independent and the change of basis is invertible. Consequently the
linear combination of rows corresponding to $bP_j-aQ_j$ cannot
annihilate all three directions when $b>0$. This nonvanishing
property survives division at each level.

\begin{proof}[Proof of Theorem~\ref{th:main}]
Set $\sigma=h-s_0^*$ and $\tau=g-h+s_0^*>0$. For
$0<\varepsilon<\tau$ and large $b$, take
$n=\lceil\log(2b)/(\tau-\varepsilon)\rceil$. Choose $j\in[n,n+2]$
with the integer $bP'_j-aQ'_j$ nonzero. The two rate estimates give
\[
 1\le |bP'_j-aQ'_j|
 \le \tfrac12+b e^{(\sigma+\varepsilon)(n+2)}|\xi-a/b|.
\]
Hence $|\xi-a/b|\gg_\varepsilon
b^{-1-(\sigma+\varepsilon)/(\tau-\varepsilon)}$ uniformly in $a$.
Letting $\varepsilon\downarrow0$ and using $\mu(\pi^2)=\mu(\xi)$
gives
\[
 \mu(\pi^2)\le 1+\frac{\sigma}{\tau}
 =\frac{g}{g-h+s_0^*}
 =\frac{821200073639400000000}{178914706731744129841}
 <\frac{4589897}{10^6}.\qedhere
\]
\end{proof}

\clearpage
\paragraph{Lean formalization.}
The \href{https://github.com/Mattie/math/blob/02d402999844a7b9524c1668d54e6314979c7569/verification/prime-window/inputs/recipes/WindowFinal.lean#L35-L48}{Lean proof of Theorem~\ref{th:main}}
uses the first-column height, error, and independence estimates of
Leibtag et al.\ \cite{L,LC} and the prime number theorem from the
PrimeNumberTheoremAnd project \cite{PNT}. The bound needs only the finite
estimate $s_{10^6}>10868226/10^8$. The proof of the full limiting rate
in~\eqref{eq:rate} is given in the text.
An \href{https://github.com/Mattie/math/actions/runs/45713249120}{independent check on GitHub Actions}
rebuilt the formalization from source. Lean's kernel and Nanoda both
checked the exact rational bound and its strict decimal consequence
in the infimum and supremum definitions of the irrationality exponent,
as well as the finiteness statement following Theorem~\ref{th:main}.

\paragraph{Acknowledgment.} OpenAI Codex agents assisted substantially
with mathematical development, Lean formalization, manuscript
preparation, and review.

\begin{thebibliography}{9}
\small\setlength{\itemsep}{3pt}
\bibitem{Casper}
R. M. Casper, \emph{The irrationality exponent of real and imaginary
quadratic periods is 2}, 9 October 2026.
\href{https://github.com/Mattie/math/blob/d9de38369ed65117d61b5230accaa9dbb930d337/preprints/The-irrationality-exponent-of-real-and-imaginary-quadratic-periods-is-2-October-9-2026/paper.pdf}{Preprint}.
\bibitem{C}
N. A. Carella, Irrationality measure of $\pi^2$, 2022.
\href{https://arxiv.org/abs/2201.00651v2}{arXiv:2201.00651v2}.
\bibitem{Z}
W. Zudilin, Two hypergeometric tales and a new irrationality measure
of $\zeta(2)$, \emph{Ann. Math. Qu\'ebec} \textbf{38} (2014), 101--117.
\href{https://doi.org/10.1007/s40316-014-0016-0}{doi:10.1007/s40316-014-0016-0}.
\bibitem{N}
D. Niedbala Giraudin, A new upper bound for the irrationality exponent
of $\zeta(2)$, 2026. \href{https://arxiv.org/abs/2610.02912}{arXiv:2610.02912}.
\bibitem{L}
E. Leibtag, H. Barkan, R. Kalisch, S. Weinbaum, M. Shalyt, and
I. Kaminer, New directions in irrationality, 2026.
\href{https://arxiv.org/abs/2610.27799v1}{arXiv:2610.27799v1}.
\bibitem{PNT}
PrimeNumberTheoremAnd contributors, \emph{PrimeNumberTheorem+
and explicit bounds}, Lean formalization.
\href{https://github.com/AlexKontorovich/PrimeNumberTheoremAnd/tree/127b50e33f6033e2147921167a117af96d9ba4ea}{Project and source}.
\bibitem{LC}
Ramanujan Machine, \emph{CMF conservation as an arithmetic tool:
a proof of $\mu(\zeta(2))<4.62675$ for the direction
$(10,9,11,-11;2,1,3)$}, supplementary derivation
for \cite{L}, 2026. \href{https://github.com/RamanujanMachine/zeta2_four_trajectory_matrices/blob/a0e8ac417cddc6ffa88e43cfd68d2b58ea339517/results/zeta2_complete_proof/zeta2_complete_proof.tex}{Version with the six-interval table and rational rate enclosure}.
\end{thebibliography}
\end{document}
